\begin{table}[htbp]
\centering
\caption{Ringkasan Historis Peta Jalan (\textit{Roadmap}) Evolusi Riset Grup Pemandu}
\label{tab:roadmap_historis}
\renewcommand{\arraystretch}{1.5}
\begin{tabular}{@{} l p{4.5cm} p{6cm} p{3.5cm} @{}}
\toprule
\textbf{Fase/Tahun} & \textbf{Fokus \& Tema Utama} & \textbf{Capaian Ilmiah Utama (\textit{Milestone})} & \textbf{Sitasi Basis} \\ \midrule

\textbf{Tahap I} & \textbf{Dinamika Fundamental Fisika Plasma} & 
Karakterisasi awal \textit{Laser-Induced Shock Wave Plasma}, eksitasi pada lingkungan tekanan rendah, dan analisis fisis deuterium/hidrogen pada padatan keras seperti pipa reaktor (\textit{Zircaloy-4}). & 
Idris et al. (2004), Idris et al. (2013) \\ \midrule

\textbf{Tahap II} & \textbf{Transisi Aplikasi: Sedimen \& Geokimia} & 
Penerapan LIBS untuk mengekstraksi dan mendeteksi kontaminasi jejak elemen pada sedimen tanah sisa Tsunami Aceh 2004 menggunakan instrumen laser TEA CO$_2$ dan Nd:YAG. & 
Idris et al. (2016), Idris et al. (2019) \\ \midrule

\textbf{Tahap III} & \textbf{Resolusi Validasi Lintas-Instrumen} & 
Analisis interdisiplin melalui silang-validasi proksi Paleo-tsunami dan geokimia tanah laut menggunakan instrumen baku tambahan seperti SEM-EDX dan mesin X-Ray Fluorescence (XRF). & 
Mitaphonna (2021), Idris et al. (2023) \\ \midrule

\textbf{Tahap IV} & \textbf{Modernisasi \textit{AI} \& Kuantifikasi Spektral} & 
\textit{(Fokus Tesis)} Mengatasi limitasi kualitatif dan \textit{matrix effect} tanah melalui pendekatan sintesis arsitektur \textit{Deep Learning} otomatis (Model regresi \textit{CNN-Transformer}). & 
Wang (2024), Walidain dkk. (2026) \\ \bottomrule
\end{tabular}
\vspace{2mm} % Space between table and source note
\par
\textit{Sumber: Tinjauan sistematis atas 100 literatur Nasrullah Idris \textit{et al.} dari Semantic Scholar (2004--2024), dirumuskan mandiri.}
\end{table}
